\documentclass[11pt,a4paper]{article}
\usepackage[utf8]{inputenc}
\usepackage[english]{babel}
\usepackage{amsmath}
\usepackage{amsfonts}
\usepackage{booktabs}
\usepackage{listings}
\usepackage{xcolor}
\usepackage{geometry}
\usepackage{hyperref}

\geometry{margin=1in}

\definecolor{codegreen}{rgb}{0,0.6,0}
\definecolor{codegray}{rgb}{0.5,0.5,0.5}
\definecolor{codepurple}{rgb}{0.58,0,0.82}
\definecolor{backcolour}{rgb}{0.95,0.95,0.92}

\lstset{
    backgroundcolor=\color{backcolour},   
    commentstyle=\color{codegreen},
    keywordstyle=\color{magenta},
    numberstyle=\tiny\color{codegray},
    stringstyle=\color{codepurple},
    basicstyle=\ttfamily\footnotesize,
    breakatwhitespace=false,         
    breaklines=true,                 
    captionpos=b,                    
    keepspaces=true,                 
    numbers=left,                    
    numbersep=5pt,                  
    showspaces=false,                
    showstringspaces=false,
    showtabs=false,                  
    tabsize=2,
    language=Python
}

\title{Parametric Automation and Optimization of a Bahtinov Focusing Mask for Catadioptric Optical Systems}
\author{Engineering Portfolio Project}
\date{\today}

\begin{document}

\maketitle

\begin{abstract}
This report outlines the programmatic generation, mechanical optimization, and structural refinement of a Bahtinov focusing mask designed specifically for the Sky-Watcher SkyMax 127T Maksutov-Cassegrain telescope. Utilizing Open-Source Computer-Aided Design (CAD) automation via FreeCAD's Python API, the geometry was mathematically derived, reinforced against structural failure with a central cross-bridge, and topologically optimized using algorithmic shape refinement techniques.
\end{abstract}

\section{Introduction and Optical Principles}
The Bahtinov mask is a specialized optical calibration device utilized in amateur astronomy and astrophotography to achieve critical focus. Operating on the principle of Fraunhofer diffraction through periodic parallel apertures (slit gratings), the mask divides the incoming wavefront into three distinct geometric clusters. 

Vertical slits yield a linear horizontal diffraction spike, while two opposing slit grids inclined at angles ($\alpha = \pm 20^\circ$) generate asymmetric diagonal spikes forming an "X" configuration. Perfect focal alignment is reached when the central horizontal spike bisects the intersection point of the diagonal spikes symmetrically.

\section{Mathematical Derivation of Parametric Properties}
For a specialized Maksutov-Cassegrain system like the SkyMax 127T, geometric calculations must integrate both the empirical diffraction grid rules and geometric space restrictions arising from central secondary mirror obstruction. The primary grid pitch ($p$) is calculated using the established empirical relationship:
\begin{equation}
p \approx \frac{F}{170}
\end{equation}
Where $F$ represents the system focal length. The slot width ($w$) is set exactly to half of the pitch to enforce a 1:1 aspect ratio, maximizing first-order diffraction intensity:
\begin{equation}
w = \frac{p}{2}
\end{equation}

\begin{table}[h]
\centering
\caption{System Constraints and Derived Model Values}
\label{tab:specs}
\begin{tabular}{lcr}
\toprule
\textbf{Parameter} & \textbf{Symbol} & \textbf{Design Value} \\
\midrule
Aperture Diameter & $D$ & 127.0 mm \\
Focal Length & $F$ & 1500.0 mm \\
Central Obstruction & $D_{obs}$ & 45.0 mm \\
Calculated Grid Pitch & $p$ & 8.8 mm \\
Resulting Slot Width & $w$ & 4.4 mm \\
Structural Bridge Width & $w_{bridge}$ & 6.0 mm \\
Target Extrusion Thickness & $t$ & 3.0 mm \\
\bottomrule
\end{tabular}
\end{table}

\section{Mechanical Optimization and Structural Integration}
A common drawback of standard algorithmically cut masks is the presence of high-aspect-ratio fragile structural elements near the tangent lines of circular intersections. To mitigate mechanical degradation and snapping risks during Fused Deposition Modeling (FDM) 3D printing, a 6.0~mm horizontal structural bridge was introduced across the absolute $X$-axis coordinate vector. 

Furthermore, sequential Boolean unions typically introduce unseemly internal seams and geometric artifacts. To address this, the topology was programmatically passed into an OpenCASCADE shape-refinement filter via Python, zipping coplannar boundaries into a single clean manifold entity.

\section{Automated Scripting Implementation}
The layout generation is controlled by a standalone Python script executed inside FreeCAD's embedded headless environment. It automatically compiles the parameters, implements boolean slicing arrays, refines internal intersection lines, and dumps output components into `.stl` and `.step` configurations.

\begin{lstlisting}[caption={Refined FreeCAD Python Automation Workflow}]
# Feature snippet highlighting programmatic topology optimization
# Complete execution workflow embedded inside bahtinov_generator.py

# ... [Geometry generation logic] ...

# Perform final Boolean subtraction from the main body
bahtinov_mask = mask_body.fuse(Part.makeCylinder(D / 2.0, thickness).cut(final_slots))

# Clean internal topological artifacts and merge coplanar segments
clean_bahtinov_mask = bahtinov_mask.refineShape()

# Save final object to Tree layout
mask_obj = doc.addObject("Part::Feature", "Bahtinov_SkyMax127_Clean")
mask_obj.Shape = clean_bahtinov_mask
doc.recompute()
\end{lstlisting}

\section{Conclusion}
The resulting workflow demonstrates a robust method for parametric CAD deployment. By replacing tedious manual drawing iterations with explicit scripting logic, modifications to underlying optical configurations can be instantly updated. The final refined asset maps directly into production pipelines for high-precision manufacturing.

\end{document}
